% Gerado por tools/tei2latex.py a partir de vol1-divisao-do-ente.xml. Regenerar, não editar.
% Recorte manual: tirada a matéria preliminar que tools/trecho.py arrastou; \pg{331} reposto.

\pg{331}

%%>>> cap-10

\chapter*{CAPITULO SEGUNDO}
\addcontentsline{toc}{chapter}{CAPITULO SEGUNDO}

\textbf{DIVISÃO DO ENTE} \textsc{SuxMARIO}: --- Ente necessario e ente contingente. --- Este infinito e ente finito.

--- Ente simples e ente composto. --- Ente immutavel e ente mutavel,

%%>>> art-12

\section*{Artigo \textbar{}.\\ Ente necessario e ente contingente.}
\addcontentsline{toc}{section}{Artigo \textbar{}. --- Ente necessario e ente contingente.}

34. Divisão do ente. --- Todo o ente --- ou é \emph{acto puro}, --- ou consta de \emph{acto} e de \emph{potencia}. Se é \emph{acto puro}, o ente é \emph{necessario}, \emph{infinito}, \emph{simples}, \emph{immutavel}. Se consta de \emph{acto} e de \emph{potencia}, é \emph{contingente}, \emph{finito}, \emph{composto}, \emph{mutavel}. D'onde a divisão de todo o. ente --- em \emph{necessario} e \emph{contingente}, --- em \emph{infinito} e \emph{finito}, --- em \emph{simples} e \emph{composto}, --- em \emph{immutavel} e \emph{mutavel}, --- Esta divisão funda-se na diversa perfeição do ente (*J). ---\footnote{Todo o ente é dotado de \emph{perfeição}. --- \emph{Perfeição é toda} a \emph{realidade}, « \emph{que} serve para fazer ou completar uma. coisa... A perfeição pode ser \emph{parcial} e \emph{fotul}. \emph{É parcial}, se representa uma das realidades, de que o ente é dotado. É \emph{total}. se representa o \emph{conjuncto} das mesmas realidades D' onde se vê que a perfeição não é um attributo especial do ente, mas é uma coisa que constitue, adequada ou inadequadamente, o proprio ente. --- Por isso um ente é e diz-se \emph{perfeito}, isto é, \emph{inteiramente feito}, emquanto possue a --- perfeição \emph{total}, de modo que não Ihe falte nenhuma das realidades. necessarias para o seu complemento. --- A' \emph{perfeição} oppüe-se a \emph{imperfeição}, que denota \emph{ausencia} do perfeição. \emph{A} imperfeição é privativa ou \emph{negativa}. \emph{É privativa}, quando é ausencia de uma perfeição \emph{devida}, e diz-se commummente \emph{privação}: tal é a ausencia de um braço no homem. É \emph{negativa}, quando é a ausencia de uma perfeição nao-devida, e chama-se \emph{negação}; tal é a ausencia das'azas no homem, ou da razão na planta. . .\par A perfeição \emph{total} divide-se em \emph{absoluta} e \emph{relativa}. \emph{A absoluta} é 0 \emph{conjuncto} de todas as realidades ou perfeições \emph{possiveis}, sem limitação alguma; tal é a perfeição de Deus, que é o proprio Ser subsistente. A \emph{relativa} éo conjuncto das realidades ou perfeigóes, que são devidas a um certo genero, ou a uma certa especie de entes: tal é a perfeição de um homem, que tem tudo o que a sua natureza exige. --- A perfeição \emph{absoluta} não admitte graus, nem divisões: a \emph{relativa} admitte graus, --- emquanto um genero ou uma especie pode abranger um maior ou menor numero de perfeições parciaes com relação a um outro genero, ou a uma outra especie, --- óu emquanto um individuo possue uma realidade ou perfeição por um modo mais \emph{infeuso que o}\pgn{332}modo por que a possue um outro individuo do mesmo genero, ou da mesma especie. --- A perfeição \emph{total} pode considerar-se --- em relação á \emph{entidade}, quando o ente tem (udo o que é indispensavel para a sua essencia e inlo- £ridade; --- em relação á \emph{operação}, quando o enté tem Ludo o que é necessario para operar em conformidade com a sua essencia; --- em relação ao \emph{fim}, quando o ente, tendo alcançado o fim pelas suas operações, encontra nelle o seu complemento, pois. o fim.é bem (Cf. Sum. \emph{Th}. p: 1, q. 6, a. 3).\par A perfeição \emph{relativa} divide-se em \emph{simples} e \emph{mista}. --- \emph{Simples} é a que no seu conceito não encerra imperfeição alguma e não se oppõe a perfeições mais elevadas; tal é a \emph{vida}, a bondade, a \emph{belleza}, etc. A perfeição \emph{simples} costuma chamar-se tambem \emph{absoluta}, mas emquanto \emph{o} seu conceito não denota imperfeição, --- \emph{Mista} é a que, embora no seu conceito exprima uma entidade. positiv. 24 ---todavia---encerra alguma. imperfeição e se oppõe a uma perfeição mais elevada: tal é a \emph{corporeidade}, que, embora denote uma realidade, comtudo importa uma imperfeição, pois o \emph{corpo é composto}, \emph{dependente}, etc., --- e se oppõe á \emph{espiritualidade}, quo é uma coisa mais elevada do que o corpo (Cf. Ansel., Monol., e. 15).\par A perfeição pode convir ao ente por tres modos: \emph{formal}, \emph{virtual} e \emph{eminente}, Convem \emph{formalmente}, quando existe no ente segundo a propria forma ou razão: especifica, isto é, segundo o-seu conceito e a sua definição; assim encontra-se o calor no \emph{fogo}, \emph{a sensibilidade} no cão, a \emph{razão} no \emph{homem}, Convem . virtualmente, quando não existe. no ente segundo a propria forma, mas existe na virtude ou causalidade do ente, emquanto este é capaz de a produzir; assim encontra-se a \emph{planta} na \emph{semente}, c todo o \emph{effeito} na sua causa. Convem eminentemente; quando existe no ente, nem segundo a propria forma, nem segundo a energia ou causalidade, mas por um modo mais nobre e mais elevado que o modo por. que se encontra nos seres, que a possuem \emph{formalmente}; assim encontram-se em Deus todas \emph{as} perfeições dos entes creados. --- Note-se que um d esses modos não exclue outro. Uma perfeição pode convir a um ente por um modo \emph{formal} e \emph{eminente} ao mesmo tempo, e convem quando existe nelle segundo a propria forma, mas pum grau mais elevado; assim a \emph{sabedoria}, que é o conhecimento das coisas pelas causas supremas, attribue-se a Deus --- não só \emph{formalmente}, emquanto, em Deus, como no homem, importa o conhecimento pelas causas supremas, --- mas lambem. \emph{eminentemente}, emquanto em \textsc{Deus} á \emph{sabedoria} é uma e a mesma coisz com essencia, ao passo que no homem é uma qualidade accidental e limitlzda. Como tambem uma perfeição pode convir a um ente por um modo \emph{viríual} e \emph{eminente} ao mesmo tempo, e convem quando não existe nelle segundo a. propria forma, mas existe na virtude ou capacidade do ente, e de um modo mais elevado e nobre; assim \emph{o raciocinio}, importando um discurso ou uma passagem da potencia para o \&cto, não se altribue a Deus \emph{formalinente}, --- mas attribue-se virtualiiente, emquanto a intelligencio \emph{divina} tem a capacidade de perceber a conclusão no principio, - c tambem \emph{eminentemente}, emquanto a intelligencia\pgn{333} divina não passa, como a do homem, do conhecido para o desconhecido, mas vé a conclusão no principio sem discurso, por intuição.\par Premissas estas noções, vejamos a razão da divisão do \textbf{ente}. A perfeição pode encontrar-se no ente --- ou de um modo \emph{necessario} Ou contingente, e temos o ente \emph{necessario} e O ente \emph{contingente}, --- ou de um modo \emph{illimitado} ou \emph{limitado}, e temos o ente \emph{infinito} e O ente \emph{finito}, --- ou de um modo simples ou \emph{composto}, e temos o ente \emph{simples} e o ente composto, --- ou de um modo \emph{permanente} ou \emph{variavel}, e lemos o ente immutavel e o ente mutavel. --- A primeira especie da divisão refere-se á \emph{existencia} do ente; a segunda, á sua \emph{coumprehkensão}; \& terceira, a sua constituição; a. quarta, a \emph{permanencia}, ou \emph{estabilidade} da perfeição. --- Como se vé, todas estas especies denotam a mesma coisa e importam os mesmos extremos, considerados sob differentes aspectos e significados por varios conceitos e termos, dos quaes um é mais claro e mais obvio que.outro: porque o Ente necessario, \emph{infinito}, simples, immutavel é Deus ; e o ente \emph{contingente}, \emph{finito}, \emph{composto}, mutavel € a \textbf{creatura}. Por isso esta divisão é a mais transcendental que se pode dar: --- Advertimos que, tratando aqui do Ente \emph{necessario}, \emph{infinito}, \emph{simples} e immutavel, não entendemos demonstrar a sua existencia; mas, admittida, pela experiencia, a existencia de entes \emph{contingentes}, \emph{finitos}, \emph{compostos} e \emph{mutaveis}, procuramos descobrir o nexo, que elles têm com o Ente necessario, \emph{infinito}, simples e \emph{immutavel}.}

